\chapter{Conclusioni}
\label{chap:conclusioni}

\section{Sintesi delle scelte architetturali}
\label{sec:sintesi}

Il progetto realizza un'istanza coerente del paradigma \emph{distributed service +
agentic operations layer} richiesto dalle linee guida. Le scelte architetturali si
caratterizzano per \emph{sobriet\`a deliberata}: meccanismi di affidabilit\`a
selezionati con motivazione esplicita, agente confinato a un ruolo di
\emph{operational assistant advisory-by-default}, MCP server intenzionalmente
read-only, soglie di confidenza e cost ceiling materializzati nel codice.

\paragraph{Reliability.}
Il sistema dimostra \textbf{graceful degradation} su tutti i piani critici: il fallimento
dell'LLM non compromette l'archiviazione; il fallimento di S3 non compromette i metadati;
il fallimento del database espone solo le operazioni che ne richiedono l'accesso. Nessuno
scenario di guasto analizzato produce perdita di dati.

\paragraph{Idempotenza.}
L'idempotenza forte sull'ingestione email (vincolo UNIQUE + watermark IMAP) garantisce
la propriet\`a \emph{at-least-once} senza duplicazione, anche in presenza di riavvii
e crash.

\paragraph{Auditabilit\`a.}
La tabella \texttt{agent\_operations} costituisce un audit trail completo e immodificabile
di ogni decisione dell'agente, incluse le operazioni di fallback. Ogni decisione \`e
tracciabile dall'email originale (conservata su S3 come \texttt{.eml}) al documento
archiviato.

\paragraph{Analisi economica.}
La trattazione SLO/costi/ROA mette in luce come il valore agentico, in questo dominio,
risieda nella riduzione del \emph{tail-risk} operativo pi\`u che nel mero risparmio
di FTE: una conclusione coerente con l'obiettivo di costruire sistemi che continuino
a ``comportarsi in modo accettabile quando le cose vanno storte''.

\section{Sviluppi futuri}
\label{sec:sviluppi_futuri}

\subsection{Deployment Kubernetes}

\TODO{I manifest Kubernetes (Deployment, Service, Ingress, StatefulSet per PostgreSQL,
HPA per backend e agent\_worker) permetteranno di:
\begin{itemize}
    \item scalare automaticamente il backend in risposta ai picchi di carico;
    \item garantire zero-downtime rolling updates;
    \item gestire i segreti tramite Kubernetes Secrets con cifratura a riposo.
\end{itemize}}

\subsection{Secondo agente: monitoring operativo}

\TODO{Il secondo agente agentico connetterà al sistema tramite un MCP server dedicato
al monitoring (Prometheus, Loki). Il suo compito sarà interpretare anomalie metriche,
correlare eventi di log e proporre azioni di remediation. Ogni proposta richiederà
approvazione dell'operatore prima dell'esecuzione; nessuna azione distruttiva sarà
automatizzata.}

\subsection{Observability stack}

\TODO{Integrazione di Prometheus per la raccolta di metriche operative, Grafana per
le dashboard, e OpenTelemetry per il tracing distribuito end-to-end attraverso tutti
i microservizi.}

\subsection{OAuth2 per account email}

\TODO{Implementazione del flusso OAuth2 completo per Google e Microsoft Outlook,
che permetterà agli utenti di collegare i propri account senza usare app password
e rispettando i requisiti di sicurezza moderni dei provider.}

\subsection{Rate limiting}

\TODO{Aggiunta del middleware \texttt{slowapi} (o configurazione a livello Ingress/API
gateway) per limitare il numero di richieste per utente e prevenire abusi degli endpoint.}

\subsection{Migration management}

\TODO{Integrazione di Alembic per la gestione delle migration incrementali dello schema
PostgreSQL, permettendo l'evoluzione del database in produzione senza perdita di dati.}
